\subsection{Module A.6: shared enclosure and local admissibility}
\label{sec:admissible-set-derivation}

\paragraph{Exact interface.}
The enclosure number $\Lambda$, canonical anchor hull $K(a,b,c)$,
admissible set $\mathcal A$, and radial envelope $c_{\max}$ are defined in
\zcref{def:finite-enclosure-number,def:finite-anchor-admissible-set}.  This
module evaluates those objects once for use by all later appendices.

\paragraph{Variables and feasible set.}
Let $\mathsf R$ be counterclockwise rotation through $2\pi/3$.  For every
nonempty compact $K\subset\mathbb R^2$ put
\[
 h_K(n):=\max_{x\in K}\langle x,n\rangle.
\]
Translate the canonical hull's vertex $V_0$ to the origin and apply a rigid
motion so that
\[
 u=\left(\frac12,\frac{\sqrt3}{2}\right),
 \qquad
 v=\left(\frac12,-\frac{\sqrt3}{2}\right).
\]
Thus its generators become
\[
 O'=0,
 \qquad A=au,
 \qquad B=bv,
 \qquad C=c(u+v),
 \qquad K=\conv\{O',A,B,C\}.
\]
For $a,b,c\in[0,1]$ set
\[
 \begin{gathered}
 s:=a+b,
 \qquad d:=a^2+ab+b^2,
 \qquad q:=s^4-s^2+ab,\\
 m:=\min\{a,b\},
 \qquad M:=\max\{a,b\}.
 \end{gathered}
\]
and
\begin{align*}
 F_L(a,b,c)&:=c^4-c^2+mc-m^2,\\
 F_T(a,b,c)&:=(s^2-1)c^2+Mc-M^2,\\
 F_S(a,b,c)&:=(m^2-1)c^2+(2mM^2+M)c+M^4-M^2.
\end{align*}

\paragraph{Objective.}
Evaluate $\Lambda(K)$ by support functions; determine the exact projection of
$\mathcal A$ to the $(a,b)$-plane; and, on that projection, calculate the
attained maximum $c_{\max}(a,b)$ with the geometric connected-component
selectors retained.

\paragraph{Exact result.}

\begin{lemma}[Shared equilateral-enclosure gauge]
\label{lem:app-shared-enclosure-gauge}
\label{prop:universal-enclosure-gauge}
\label{prop:new-enclosure-gauge}
For every nonempty compact $K$,
\begin{equation}
 \Lambda(K)=\frac2{\sqrt3}
 \min_{\lVert n\rVert=1}
 \sum_{j=0}^2h_K(\mathsf R^jn).
 \label{eq:app-enclosure-gauge}
\end{equation}
If $K$ is a polygon with at least two distinct points, a minimizing normal
may be chosen perpendicular to a hull edge.
\end{lemma}

\begin{lemma}[Exact local admissible set and radial envelope]
\label{lem:app-local-admissible-set}
\label{prop:exact-admissible-set}
\label{prop:new-exact-local-set}
The exact projection is
\begin{equation}
 \{(a,b)\in[0,1]^2:\exists c\in[0,1]\ (a,b,c)\in\mathcal A\}
 =\{(a,b)\in[0,1]^2:a^2+ab+b^2\le1\}.
 \label{eq:app-admissible-projection}
\end{equation}
On $[0,1]^3$, the admissible set is
\begin{equation}
 \mathcal A=\mathcal A_L\cup\mathcal A_T\cup\mathcal A_S,
 \label{eq:derived-admissible-cells}
\end{equation}
where
\begin{align*}
 \mathcal A_L&=\{d\le1,\ s\le1,\ q\le0,\ F_L\le0\},\\
 \mathcal A_T&=\{d\le1,\ s\le1,\ q\ge0,\ F_T\le0,\ c\le2M\},\\
 \mathcal A_S&=\{d\le1,\ s>1,\ F_S\le0,\ c\le1/2\}.
\end{align*}
It is compact, coordinatewise downward closed, and invariant under exchange
of $a,b$.  Every $c$-fiber therefore has an attained maximum, given by
\begin{equation}
 c_{\max}(a,b)=
 \begin{cases}
  1,&a=b=0,\\
  C_L(m),&s\le1,\ q\le0,\ (a,b)\ne(0,0),\\
  \dfrac{2M}{1+\sqrt{4s^2-3}},&s\le1,\ q>0,\\[3mm]
  \dfrac{M(2mM+1)-M\sqrt{4d-3}}{2(1-m^2)},&s>1,
 \end{cases}
 \label{eq:radial-envelope}
\end{equation}
where $C_L(m)$ is the unique root in $[\sqrt3/2,1]$ of
\[
 z^4-z^2+mz-m^2=0.
\]
In particular, $c_{\max}(a,b)=c_{\max}(b,a)$.
\end{lemma}

\paragraph{Solution.}

\begin{proof}[Proof of \zcref{lem:app-shared-enclosure-gauge}]
Fix outward unit normals $n,\mathsf Rn,\mathsf R^2n$.  Translating the three
supporting lines to the support numbers of $K$ gives the least containing
triangle with those normals.  Its altitude is the sum of the three support
numbers, because $n+\mathsf Rn+\mathsf R^2n=0$, and its side is that sum
divided by $\sqrt3/2$.  Minimizing over the unit circle gives
\eqref{eq:app-enclosure-gauge}; continuity of the support function gives the
minimum.

For a polygon, partition the unit circle into angular cells on which the
three support vertices are fixed.  On each open cell the support sum has the
form $A\cos\theta+B\sin\theta$ and its second derivative is its negative.
An interior critical point is therefore a maximum unless the expression is
constant, so a minimum occurs at a cell endpoint, where a support tie makes
the normal perpendicular to a hull edge.
\end{proof}

\begin{proof}[Proof of \zcref{lem:app-local-admissible-set}]
Assume first $s>0$ and $cs\le ab$.  If $a,b>0$, then
\[
 \frac ca+\frac cb\le1,
\]
so
\[
 C=\frac caA+\frac cbB+\left(1-\frac ca-\frac cb\right)O'
 \in\conv\{O',A,B\}.
\]
The zero-coordinate cases follow directly.  The distance $AB$ is $\sqrt d$.
With
\[
 T:=-bu-av,
\]
the triangle $\conv\{A,B,T\}$ is equilateral of side $\sqrt d$, and
\[
 O'=\frac{b^2}{d}A+\frac{a^2}{d}B+\frac{ab}{d}T.
\]
Thus it contains $K$, while every enclosing triangle has diameter at least
$AB$.  Hence $\Lambda(K)=\sqrt d$.  When $a=b=0$, the hull is the segment
$[O',C]$ and its enclosure infimum is $c$.

Now suppose $s>0$ and $cs>ab$.  The hull vertices have cyclic order
$O',B,C,A$.  By \zcref{lem:app-shared-enclosure-gauge}, it suffices to test
the four edge normals.  For $O'A$ the three supports give
\[
 L_{O'A}=b+\max\{a,c\},
\]
and reflection gives $L_{BO'}=a+\max\{b,c\}$.  For $AC$, put
\[
 R_a:=\sqrt{a^2-ac+c^2},
 \qquad
 n_0:=\frac1{2R_a}(\sqrt3a,2c-a).
\]
The supports in the directions $n_0,\mathsf Rn_0,\mathsf R^2n_0$ are
\[
 \frac{\sqrt3ac}{2R_a},
 \qquad
 \frac{\sqrt3a(a-c)_+}{2R_a},
 \qquad
 \frac{\sqrt3c\max\{b,(c-a)_+\}}{2R_a}.
\]
Consequently
\[
 L_{AC}=
 \begin{cases}
 (a^2+bc)/R_a,&a\ge c,\\
 c(a+b)/R_a,&a<c\le a+b,\\
 c^2/R_a,&c>a+b,
 \end{cases}
\]
and $L_{CB}$ is obtained by interchanging $a,b$.  Therefore
\begin{equation}
 \Lambda(K)=\min\{L_{O'A},L_{BO'},L_{AC},L_{CB}\}.
 \label{eq:app-four-local-calipers}
\end{equation}

Assume by symmetry that $a=m\le b=M$.  Reading the active supports in
\eqref{eq:app-four-local-calipers} gives
\[
\begin{array}{c|c|c}
\text{range}&\text{active supports}&\text{unit-frontier equation}\\ \hline
s<1,\ q<0&\{O'\},\{C\},\{A,C\}&F_L=0\\
s<1,\ q>0&\{O'\},\{B,C\},\{A\}&F_T=0\\
s=1&\text{preceding pattern with an $O',B$ tie}&c=M\\
s>1&\{B\},\{B,C\},\{A\}&F_S=0.
\end{array}
\]
The inactive-support inequalities reduce, respectively, to $q\le0$,
$q\ge0$, and $s>1$.  Squaring the positive side-length equations gives
exactly $F_L=0,F_T=0,F_S=0$.

For the $T$ cell the roots are
\[
 c_- =\frac{2M}{1+\sqrt{4s^2-3}},
 \qquad
 c_+ =\frac{2M}{1-\sqrt{4s^2-3}}.
\]
The squared inequality contains $0\le c\le c_-$ and a formal high component
$c\ge c_+$.  Since
\[
 c_-\le2M\le c_+,
\]
with equality throughout only at root coalescence, $c\le2M$ selects exactly
the component connected to $c=0$.

For the $S$ cell,
\[
 F_S(a,b,0)=M^2(M^2-1)<0,
 \qquad
 F_S(a,b,M)=M^2(s^2-1)>0.
\]
Moreover, $4F_S(a,b,1/2)$ is increasing in $m$, and at $m=1-M$ equals
$M^2(2M-1)^2$.  Since $s>1$ gives $m>1-M$, one has
$F_S(a,b,1/2)>0$.  The quadratic opens downward, so $1/2$ lies strictly
between its roots and $c\le1/2$ selects the smaller-root component containing
$c=0$.

For the $L$ cell, the fiber issuing from $c=0$ reaches one selected root in
$[\sqrt3/2,1]$; uniqueness there follows from
$4c^3-2c+m>0$.  At $q=0$ the $L$ and $T$ formulas meet and both give $c=s$.
This proves the three exact cells and solving their selected frontier
equations gives \eqref{eq:radial-envelope}.

At $c=0$, the preceding triangular calculation shows admissibility exactly
when $d\le1$; if $d>1$, the two anchors $A,B$ are more than unit distance
apart.  This proves \eqref{eq:app-admissible-projection}.  The support formula
makes $\Lambda(K(a,b,c))$ continuous, hence $\mathcal A$ is compact.
Reflection exchanges $a,b$.  Reducing any coordinate moves its anchor toward
the already present point $O'$, so the new hull lies in the old hull; this
proves coordinatewise downward closure.  Its nonempty $c$-fibers are compact
initial intervals, proving attainment of $c_{\max}$.
\end{proof}

\paragraph{Body consequence.}
The support formula is the Appendix A supplier for
\zcref{def:finite-enclosure-number}.  The projection, component selectors,
attainment, symmetry, and formula \eqref{eq:radial-envelope} supply
\zcref{def:finite-anchor-admissible-set} and every later use of
$c_{\max}$.
